\subsection{DET-15(d): actual local instantiation and weighting restrictions}

The local hypotheses of \fid{DET-15} are instantiated in the source's
final period construction.  The relevant assertion retains the source's
standing global data and numerical regime; it does not assert the
unconditioned existence of an arbitrary \texttt{D7} or
\texttt{HidaSetting}.

The actual construction is
\path{Development/AbelianVarietyTheory/Proofs/DecoratedLatticeProofs/ExistsTermLoopOfStandingHypotheses.lean},
lines 59--119.  From its displayed \texttt{HidaSetting}, \texttt{Full},
\texttt{FinX}, \texttt{SIso}, \texttt{Ties}, standing hypotheses and
\texttt{ConcTie}, it chooses \texttt{Choice X 1} at line 66, obtains
\texttt{WuBessel} at lines 67--69, and constructs
\texttt{GoodW} and \texttt{GoodS} at lines 70--74.  These are supplied
to the actual local and loop objects at lines 104--113.
They are constructed here, rather than assumed as extra free predicates.

The measures and vectors in this construction are specified by
\path{Formalism/LocalToGlobal/SiegelLatticeTheory/FlagStabilizerTensors/LocWSOf.lean}.
Its \texttt{Choice}, lines 78--91, includes transport matrices, a
translated lowest-type vector, normalized compact Haar measures, and
quotient measures with their \texttt{IsQuotMeas} identities.
The actual existence proof in
\path{Formalism/LocalToGlobal/SiegelLatticeTheory/FlagStabilizerTensors/RootLatticeThresholds/Choice.lean},
lines 574--611, chooses the compact Haar measures at lines 596--603,
the quotient measures at lines 604--610, and constructs
\texttt{Choice.mk} at line 611.  In \path{LocWSOf.lean},
lines 109--121 define \texttt{locW}; lines 138--145 prove
\texttt{GoodW}; lines 153--175 choose the actual Bessel maps and
their quotient-norm normalization; lines 200--207 define
\texttt{locS}.  The condition \texttt{GoodSRest}, lines 212--223,
contains a smooth nonzero base vector, a strictly positive quotient
normalization, both finite quotient-norm identities, and the formal
degree.  Its final proof is
\path{Formalism/LocalToGlobal/SiegelLatticeTheory/FlagStabilizerTensors/RootLatticeThresholds/IsLowUnitHBaseVector.lean},
lines 23--27; the immediate positive-norm constructor is
\path{Formalism/LocalToGlobal/SiegelLatticeTheory/FlagStabilizerTensors/RootLatticeThresholds/ChoiceFinal.lean},
lines 17--22.

Unless a different directory is displayed, abbreviated analytic proof
filenames below are in
\path{Development/AbelianVarietyTheory/Proofs/FlagMeasures/}.
The predicates \texttt{LocIn} and \texttt{BSch} in
\path{Development/AbelianVarietyTheory/Definitions/SemilocalFrames/DegVand.lean},
lines 51--63, supply exactly the multiplicity-one, orbit-span decay,
and quotient-norm scaling needed by the \texttt{IsEll} constructor.
\path{AmplitudeMultOneDecayOfLargeGapHypotheses.lean}, lines 45--78,
supplies the first two; \path{SchurAlongFamily.lean}, lines 26--33,
supplies \texttt{BSch}.  The final threshold assembly is
\path{Development/AbelianVarietyTheory/Proofs/DecoratedLatticeProofs/SchurAtRecordGap.lean},
lines 54--70.  At line 56 it obtains the source gap input from
\texttt{uniqGap\_and\_besselOfW}; at lines 68--70 it supplies the
actual \texttt{BSch} and \texttt{LocIn} inputs to the Schur estimate.
The former combined source result is
\path{BesselMultiplicityOneLargeGapAndWhittakerBessel.lean}, lines 9--10.
Its named imported results are retained as source inputs.

For the resulting local data \(W,A\), put
\[
 k=D.R.k,\quad r=D.R.rS,\quad n=2r+k+1,\quad
 \sigma=A.\sigma,\quad w_0=\frac{A.v_0}{\|A.v_0\|},\quad a=2n^5.
\]
The group measure is \(W.\mu_G\), the subgroup measure is
\(\mu_S=\operatorname{muProd}(W.\mu_W,W.\mu_N)\), and the quotient
measure is \(W.\nu_S\).  The side objects are \(z=W.z\) and
\(\pi=W.\pi\), with the constant character parameters
\[
 c_j=\frac{2D.\mathrm{host}.\mathrm{ka}}
 {D.\mathrm{host}.g(0)\,D.\mathrm{host}.g(1)}.
\]
The positive Haar density \(c_H\) is chosen in
\path{MeasurableBcWeight.lean}, lines 208--209, and its Schur
normalization is
\[
 d=\frac{A.\mathrm{deg}}{c_H\,\operatorname{volUn}(n)}>0.
\]
That file constructs both \texttt{SideIrr} assertions at lines
203--206, obtains the quotient identities at lines 198--199,
and constructs precisely the normalized \texttt{StdSm} at lines
222--226.  It supplies the character conversion and the two norm
scaling identities at lines 228--234, then invokes the exported
amplitude theorem at lines 261--263.

The proof of that theorem,
\path{LintegralHilbertSchmidtNormSqBesselAmplitudeLe.lean},
lines 44--119, constructs \(\ell\) with the exact
\texttt{IsEll} hypothesis of \fid{DET-15}.
The nonzero square-integrable period gives \texttt{NonVan} at
lines 73--83; line 84 invokes
\texttt{exists\_besselMap\_rankOne\_factorisation\_translateSpan}.
Its scalar normalization is
\[
 B(\sigma(g)A.v_0)=t\,\ell(\sigma(g)w_0),\qquad
 |t|^2=A.cN\,\dim W_p\,d.
\]
Line 105 actually invokes
\texttt{besselModel\_projector\_on\_line} to establish this
normalization.  Lines 124--126 construct these functionals on
both sides of the amplitude.  This is a concrete source
instantiation of the local operator model under the final standing
hypotheses.  It supplies no additional identification of its
spectral weights with the zeros of \(\Xi\).

\paragraph{Which displayed weights are available.}
Several nearby uses of \emph{weighted} have a different, explicit
mathematical type.  The complete proofs
\path{CompanionTheory/AutomorphicTheory/PieriRules/ThetaBesselIsometries/TensorWeightedBesselSum.lean},
lines 21--62, and the adjacent
\path{TensorWeightedBesselGap.lean}, lines 65--81, concern finite
\texttt{Fintype}-indexed decompositions with coefficients
\[
 \epsilon_I\kappa_{1,I}\overline{\kappa_{0,I}}s
 \quad\hbox{or}\quad
 \epsilon_I\kappa_{1,I}\overline{\kappa_{0,I}}s/s_I.
\]
The \texttt{TensH} definition in
\path{Development/AbelianVarietyTheory/Definitions/Jets/SumIoPrAlt.lean},
lines 88--94, explicitly allows \(\epsilon_I=\pm1\).
The exact finite period coefficient, including the formal degree,
is exported in
\path{Development/AbelianVarietyTheory/Proofs/FlagMeasures/BlockProjectionBesselStatement.lean},
lines 9--10.  The adjacent coefficient construction in
\path{CompanionTheory/AutomorphicTheory/PieriRules/ThetaBesselIsometries/FamilyFunctionalCoefficient.lean},
lines 21--62, has the fixed cutoff \(n_1(\nu)\le q_0\), coefficient
norm \(|s_\nu|=d_\nu\) on its nonzero terms, and zero functions
outside that cutoff.  None of these displayed identities states
that its coefficients are nonnegative summable eigenvalues on
the vector index of \fid{DET-15}.

Likewise,
\path{ReciprocityTheory/ReductiveGroups/DiscreteSeries/FormalDegree/CoefLoc.lean},
lines 11--27, makes \texttt{CoefWeightedLoc} a spatial integral
of a matrix coefficient times a matrix-norm power.
\path{CompanionTheory/AutomorphicTheory/BesselWaveGrids/WeightedUniformRotatedOfExpansion.lean},
lines 16--27, places a power of a finite minor in an orbital-norm
inequality.  Those specific interfaces supply spatial estimates,
not the missing spectral identity.  This conclusion is limited
to the displayed interfaces and uses read here.

\paragraph{Constant, idempotent, and invariant weighting obstructions.}
For the actual orthonormal model vectors \(\eta_j\) of
\fid{DET-15}, suppose \(T\eta_j=\beta_j\eta_j\) and
\(|\beta_j|\ge\varepsilon>0\) for infinitely many \(j\).
For two such distinct indices,
\[
 \|T\eta_i-T\eta_j\|^2=|\beta_i|^2+|\beta_j|^2
 \ge2\varepsilon^2.
\]
The image of the unit ball is not relatively compact, so \(T\)
cannot be compact or trace class.  This rules out an infinite
\(0/1\) choice and any fixed nonzero scalar weight.
Unweighted partial sums with \(N\) selected terms have trace \(N\);
distinct nested sums have operator-norm difference one and
trace-norm difference equal to the number of new terms.
Finite \(0/1\) choices yield only determinant polynomials.
The source formal degree is constant within a fixed representation
and Haar measure, as proved in \fid{DET-15}; weighting every
vector by \(d\) or \(d^{-1}\) therefore does not supply decay.

There is also an obstruction if full representation invariance is
required.  A nonzero positive compact \(C\) commuting with every
\(\sigma(g)\) on the original irreducible infinite-dimensional
\(H_\sigma\) would have a top positive eigenvalue with a nonzero
finite-dimensional eigenspace.  Commutation makes this closed
space invariant, and irreducibility makes it all \(H_\sigma\),
a contradiction.  Hence such a \(C\) is zero.
This argument imposes commutation only on the proposed operator
to which it is applied; it does not assert that all source
weights, or the arbitrary weights of \fid{DET-15}, commute with
the full representation.

\subsection{DET-17: prime weights give an arithmetic trace, but a different determinant}

The source does supply useful nonnegative summable arithmetic
coefficients.  There is a precise affirmative construction using
them in the actual common Bessel model.  Its identification with
the model projections is a new combination made here.

\begin{proposition}[DET-17]
Retain one local source model and the actual orthogonal rank-one
convolution projections \(P_n=\eta_n\eta_n^*\) of \fid{DET-15},
indexed by \(n\ge2\).  For \(\Re s>1\), define
\begin{equation}\label{eq:bessel-prime-operator}
 A(s)=\sum_{n\ge2}\Lambda(n)n^{-s}P_n,\qquad
 A_N(s)=\sum_{2\le n\le N}\Lambda(n)n^{-s}P_n .
\end{equation}
Then \(A(s)\) is normal and trace class, with
\begin{align}
 \|A(s)-A_N(s)\|_1
 &=\sum_{n>N}\Lambda(n)n^{-\Re s},
 \label{eq:bessel-prime-tail}\\
 \Tr A(s)&=-\frac{\zeta'(s)}{\zeta(s)},\\
 D_s(w):=\det(I-wA(s))
 &=\prod_{n\ge2}\bigl(1-w\Lambda(n)n^{-s}\bigr).
 \label{eq:bessel-prime-det}
\end{align}
For real \(\sigma>1\), \(A(\sigma)\) and all its finite
approximants are positive.  Each finite approximant is the
bounded extension of an actual source convolution.

At \(\sigma=2\), \(D_2\) is not \(G(w)=\Xi(\sqrt w)/\Xi(0)\),
even after any fixed nonzero rescaling of \(w\) or multiplication
by a zero-free entire factor.  This obstruction requires neither
RH nor the prime number theorem.
\end{proposition}

\begin{proof}
The source prime coefficients and their normalization are exact.
\path{Formalism/PeriodIntegrals/UnitaryGroupMassFormula/ExplicitFormula/PerronTruncOf.lean},
lines 45--52, constructs their summable positive real majorant,
using the general von Mangoldt Dirichlet-series theorem; lines
53--60 bound complex terms by that majorant.
\path{Formalism/ZeroFreeRegion/ZetaLogDerivativeBounds/NormLogDerivLeLamF.lean},
lines 10--31, gives the termwise identification with
\(\Lambda(n)n^{-\sigma}\), positivity, absolute summability,
and the log-derivative identity.  The latter is also used
explicitly in
\path{Development/AutomorphicArithmetic/DedekindLFunctions/LandauTheorems/mellinGchiModOne.lean},
lines 168--193, after the absolute-convergence argument for the
actual prime-ramp Mellin transform.

For completeness, summability also follows from
\(0\le\Lambda(n)\le\log n\), used explicitly in
\path{Formalism/PeriodIntegrals/UnitaryGroupMassFormula/ExplicitFormula/PerronSumBound.lean},
lines 173 and 190--195.  If \(\sigma>1\), choose
\(\epsilon=(\sigma-1)/2>0\).  The elementary bound
\(\log x\le x^\epsilon/\epsilon\), \(x\ge1\), gives
\[
 \Lambda(n)n^{-\sigma}
 \le\epsilon^{-1}n^{-(\sigma+1)/2},
\]
whose sum converges.  For complex \(s\), the absolute value of
the coefficient in \eqref{eq:bessel-prime-operator} is exactly
\(\Lambda(n)n^{-\Re s}\).

On the orthonormal family \(\eta_n\), the operator in
\eqref{eq:bessel-prime-operator} is diagonal with those
coefficients, and is zero on the orthogonal complement of
their span.  Its singular values are their absolute values.
This proves trace class, normality, and
\eqref{eq:bessel-prime-tail}.  The continuous trace on the
trace ideal gives
\[
 \Tr A(s)=\sum_{n\ge2}\Lambda(n)n^{-s}
 =-\zeta'(s)/\zeta(s),\qquad \Re s>1,
\]
with the sign and normalization of the cited source identity.
For real \(\sigma>1\) the diagonal coefficients are
nonnegative, proving positivity.

The finite operators are actual source convolutions with
\[
 e_{N,s}(g)=
 \sum_{2\le n\le N}\Lambda(n)n^{-s}
 \,\mathrm{eFun}(\mu_G,\sigma_{\rm rep},u_n)(g).
\]
There is no further formal-degree factor: each source
\(\mathrm{eFun}\) already gives \(P_n\), whose trace is one.
The infinite operator also has the actual integral kernel
\[
 K_s(x,y)=\sum_{n\ge2}\Lambda(n)n^{-s}
                 \eta_n(x)\eta_n(y)^*
\]
with the almost-everywhere absolute convergence and product
\(L^2\) convergence proved in \fid{DET-15}.
No unproved infinite group-convolution kernel is asserted.
The ordinary trace-class determinant formula gives
\eqref{eq:bessel-prime-det}.  Absolute summability makes this
product locally uniformly convergent in \(w\).
Away from the zeros of its individual nonzero factors,
the tail has an absolutely convergent logarithm series;
therefore these factors give all its zeros, with their full
coincident-factor multiplicities.

At \(s=2\), its zeros are consequently
\[
 w_n=\frac{n^2}{\Lambda(n)}
 \quad\text{for }\Lambda(n)>0.
\]
Let \(N_2(R)\) count them in \(|w|\le R\), including
multiplicity.  For each counted \(n\ge2\),
\[
 n^2\le R\Lambda(n)\le R\log n<Rn,
\]
so \(n<R\), and then \(n^2\le R\log R\).  Counting all integers
rather than just prime powers gives, for sufficiently large \(R\),
\begin{equation}\label{eq:bessel-prime-zerocount}
 N_2(R)\le\lfloor\sqrt{R\log R}\rfloor.
\end{equation}
This elementary upper bound needs no information about the
distribution of primes.

In contrast, the classical zero count cited in \fid{B-01},
\cite[Theorem 9.4, (9.4.3), p.~214]{titch}, is
\[
 N_\zeta(T)=\frac{T}{2\pi}\log\frac{T}{2\pi}
            -\frac{T}{2\pi}+O(\log T).
\]
Write \(\gamma=-i(\rho-\tfrac12)\).  The functional equation
pairs \(\rho\) with \(1-\rho\), with equal multiplicity, so each
nonreal pair contributes once to the zeros \(w=\gamma^2\) of
\(G\).  The critical-strip bound gives
\[
 (\Im\rho)^2\le|\rho-\tfrac12|^2
             \le(\Im\rho)^2+\tfrac14.
\]
Allowing an irrelevant bounded number of possible real zeros,
the number \(N_G(R)\) of zeros of \(G\) in \(|w|\le R\) satisfies
\[
 N_\zeta(\sqrt{R-\tfrac14})\le N_G(R)
 \le N_\zeta(\sqrt R)+O(1).
\]
The two sides have the same leading term; hence
\begin{equation}\label{eq:bessel-Xi-zerocount}
 N_G(R)\sim\frac{1}{4\pi}\sqrt R\log R.
\end{equation}
The upper bound \eqref{eq:bessel-prime-zerocount} is
\(o(\sqrt R\log R)\), so the two determinants cannot agree.
For fixed \(a\ne0\), the zero count of \(D_2(aw)\) is
\(N_2(|a|R)=O_a(\sqrt{R\log R})\), giving the same contradiction.
A zero-free entire multiplier does not change the zero count.
\end{proof}

Thus positive arithmetic weights and the exact trace
\(-\zeta'/\zeta\) are available on an actual source period space.
The trace identity concerns the parameter \(s\), while the
Fredholm determinant concerns \(w\).  It does not identify
\(A(s)\) with the resolvent of the fixed Xi operator in
\fid{DET-11}.  At \(s=2\) the explicit zero-count obstruction
proves that even this natural choice fails the required
determinant identity.  The scalar-endpoint reconstruction in
\fid{DET-16} supplies a different arithmetic prescription under
RH. The remaining direct-sign question is the identification of
the historical finite source family with that prescription and
its positivity without passing through the RH-equivalent kernel
criterion.
